\appendix

\section{Reviewed Corpus}
\label{app:corpus}

Tables~\ref{tab:reviewed-corpus-1}--\ref{tab:reviewed-corpus-3} list the 52 papers included in the review.

\begin{table*}[p]
\centering
\begin{tabular}{p{0.8cm}p{11.1cm}p{3.0cm}}
\toprule
\textbf{No.} & \textbf{Citation and title} & \textbf{Venue} \\
\midrule
1 &
\citealp{Xiao2024HealMe}. \textit{HealMe: Harnessing Cognitive Reframing in Large Language Models for Psychotherapy} &
ACL \\
\midrule
2 &
\citealp{Gabriel2024Can}. \textit{Can AI Relate: Testing Large Language Model Response for Mental Health Support} &
EMNLP \\
\midrule
3 &
\citealp{Haydarov2025Towards}. \textit{Towards AI-Assisted Psychotherapy: Emotion-Guided Generative Interventions} &
EMNLP \\
\midrule
4 &
\citealp{Qiu2025EmoAgent}. \textit{EmoAgent: Assessing and Safeguarding Human-AI Interaction for Mental Health Safety} &
EMNLP \\
\midrule
5 &
\citealp{Aghakhani2025From}. \textit{From Conversation to Automation: Leveraging LLMs for Problem-Solving Therapy Analysis} &
ACL \\
\midrule
6 &
\citealp{Wang2025AnnaAgent}. \textit{AnnaAgent: Dynamic Evolution Agent System with Multi-Session Memory for Realistic Seeker Simulation} &
ACL \\
\midrule
7 &
\citealp{Wang2025CARE-Bench}. \textit{CARE-Bench: A Benchmark of Diverse Client Simulations Guided by Expert Principles for Evaluating LLMs in Psychological Counseling} &
AAAI \\
\midrule
8 &
\citealp{Sungarda2025GenAI}. \textit{GenAI for Social Work Field Education: Client Simulation with Real-Time Feedback} &
IEEE BigData \\
\midrule
9 &
\citealp{Liu2023ChatCounselor}. \textit{ChatCounselor: A Large Language Models for Mental Health Support} &
CIKM 2023 \\
\midrule
10 &
\citealp{Chiu2024A}. \textit{A Computational Framework for Behavioral Assessment of LLM Therapists} &
arXiv \\
\midrule
11 &
\citealp{Li2024Understanding}. \textit{Understanding the Therapeutic Relationship between Counselors and Clients in Online Text-based Counseling using LLMs} &
arXiv \\
\midrule
12 &
\citealp{Wang2024Towards}. \textit{Towards a Client-Centered Assessment of LLM Therapists by Client Simulation} &
arXiv \\
\midrule
13 &
\citealp{Hu2024PsycoLLM}. \textit{PsycoLLM: Enhancing LLM for Psychological Understanding and Evaluation} &
IEEE \\
\midrule
14 &
\citealp{Fouda2026PsychiatryBench}. \textit{PsychiatryBench: a multi-task benchmark for LLMs in psychiatry} &
npj Digital Medicine \\
\midrule
15 &
\citealp{Cai2026Exploring}. \textit{Exploring Safety Alignment Evaluation of LLMs in Chinese Mental Health Dialogues via LLM-as-Judge} &
arXiv \\
\midrule
16 &
\citealp{Qiu2026PsyCLIENT}. \textit{PsyCLIENT: Client Simulation via Conversational Trajectory Modeling for Trainee Practice and Model Evaluation in Mental Health Counseling} &
arXiv \\
\midrule
17 &
\citealp{Nguyen2025Do}. \textit{Do Large Language Models Align with Core Mental Health Counseling Competencies?} &
NAACL 2025 \\
\midrule
18 &
\citealp{Zhang2025CBT-BENCH}. \textit{CBT-BENCH: Evaluating Large Language Models on Assisting Cognitive Behavior Therapy} &
ACL \\
\bottomrule
\end{tabular}
\caption{Papers included in the 52-paper review corpus (Part 1 of 3).}
\label{tab:reviewed-corpus-1}
\end{table*}

\begin{table*}[p]
\centering
\begin{tabular}{p{0.8cm}p{11.1cm}p{3.0cm}}
\toprule
\textbf{No.} & \textbf{Citation and title} & \textbf{Venue} \\
\midrule
19 &
\citealp{Kim2025KMI}. \textit{KMI: A Dataset of Korean Motivational Interviewing Dialogues for Psychotherapy} &
ACL \\
\midrule
20 &
\citealp{Deshpande2025Evaluating}. \textit{Evaluating Large Language Models for Enhancing Live Chat Therapy: A Comparative Study with Psychotherapists} &
ACL \\
\midrule
21 &
\citealp{Jaiswal2024Building}. \textit{Building Personality-Adaptive Conversational AI for Mental Health Therapy} &
ACM BCB '24 \\
\midrule
22 &
\citealp{Xu2025MentalChat16K}. \textit{MentalChat16K: A Benchmark Dataset for Conversational Mental Health Assistance} &
ACM KDD‘25 \\
\midrule
23 &
\citealp{Schmidmaier2025Using}. \textit{Using a Secondary Channel to Display the Internal Empathic Resonance of LLM-Driven Agents for Mental Health Support} &
ACM ICMI-MLMI'25 \\
\midrule
24 &
\citealp{Schmidmaier2025UsingNonverbal}. \textit{Using Nonverbal Cues in Empathic Multi-Modal LLM-Driven Chatbots for Mental Health Support} &
ACM \\
\midrule
25 &
\citealp{Ozgun2025Trustworthy}. \textit{Trustworthy AI Psychotherapy: Multi-Agent LLM Workflow for Counseling and Explainable Mental Disorder DiagNosis} &
ACM CIKM '25 \\
\midrule
26 &
\citealp{Louie2026Can}. \textit{Can LLM-Simulated Practice and Feedback Upskill Human Counselors? A Randomized Study with 90+ Novice Counselors} &
ACM CHI '26 \\
\midrule
27 &
\citealp{Kuhail2025Human-Human}. \textit{Human-Human vs Human-AI Therapy: An Empirical Study} &
IJHCI \\
\midrule
28 &
\citealp{Louie2024Roleplay-doh}. \textit{Roleplay-doh: Enabling Domain-Experts to Create LLM-simulated Patients via Eliciting and Adhering to Principles} &
EMNLP \\
\midrule
29 &
\citealp{Wang2024PATIENT}. \textit{PATIENT-$\Psi$: Using Large Language Models to Simulate Patients for Training Mental Health Professionals} &
EMNLP \\
\midrule
30 &
\citealp{Lee2024CACTUS}. \textit{CACTUS: Towards Psychological Counseling Conversations using Cognitive Behavioral Theory} &
EMNLP \\
\midrule
31 &
\citealp{Basar2024To}. \textit{To What Extent Are Large Language Models Capable of Generating Substantial Reflections for Motivational Interviewing Counseling Chatbots?A Human Evaluation} &
ACL \\
\midrule
32 &
\citealp{Yang2025Consistent}. \textit{Consistent Client Simulation for Motivational Interviewing-based Counseling} &
ACL \\
\midrule
33 &
\citealp{Yang2025CAMI}. \textit{CAMI: A Counselor Agent Supporting Motivational Interviewing through State Inference and Topic Exploration} &
ACL \\
\midrule
34 &
\citealp{Xie2025PsyDT}. \textit{PsyDT: Using LLMs to Construct the Digital Twin of Psychological Counselor with Personalized Counseling Style for Psychological Counseling} &
ACL \\
\midrule
35 &
\citealp{Kermani2025Systematic}. \textit{A Systematic Evaluation of LLM Strategies for Mental Health Text Analysis: Fine-tuning vs. Prompt Engineering vs. RAG} &
ACL CLPsych 2025 \\
\bottomrule
\end{tabular}
\caption{Papers included in the 52-paper review corpus (Part 2 of 3).}
\label{tab:reviewed-corpus-2}
\end{table*}

\begin{table*}[p]
\centering
\begin{tabular}{p{0.8cm}p{11.1cm}p{3.0cm}}
\toprule
\textbf{No.} & \textbf{Citation and title} & \textbf{Venue} \\
\midrule
36 &
\citealp{Basar2025How}. \textit{How Well Can Large Language Models Reflect? A Human Evaluation of LLM-generated Reflections for Motivational Interviewing Dialogues} &
ACL COLING \\
\midrule
37 &
\citealp{Feng2025Reframe}. \textit{Reframe Your Life Story: Interactive Narrative Therapist and InNovative Moment Assessment with Large Language Models} &
EMNLP 2025 \\
\midrule
38 &
\citealp{Xu2025MultiAgentESC}. \textit{MultiAgentESC: A LLM-based Multi-Agent Collaboration Framework for Emotional Support Conversation} &
EMNLP 2025 \\
\midrule
39 &
\citealp{Kim2025MIRROR}. \textit{MIRROR: Multimodal Cognitive Reframing Therapy for Rolling with Resistance} &
EMNLP 2025 \\
\midrule
40 &
\citealp{Kim2025Can}. \textit{Can You Share Your Story? Modeling Clients' Metacognition and Openness for LLM Therapist Evaluation} &
ACL 2025 \\
\midrule
41 &
\citealp{Liu2025Eeyore}. \textit{Eeyore: Realistic Depression Simulation via Expert-in-the-Loop Supervised and Preference Optimization} &
ACL \\
\midrule
42 &
\citealp{Wang2025Feel}. \textit{Feel the Difference? A Comparative Analysis of Emotional Arcs in Real and LLM-Generated CBT Sessions} &
EMNLP \\
\midrule
43 &
\citealp{Samuel2025PersonaGym}. \textit{PersonaGym: Evaluating Persona Agents and LLMs} &
EMNLP \\
\midrule
44 &
\citealp{Chen2025MIND}. \textit{MIND: Towards Immersive Psychological Healing with Multi-Agent Inner Dialogue} &
EMNLP \\
\midrule
45 &
\citealp{Berrezueta-Guzman2024Future}. \textit{Future of ADHD Care: Evaluating the Efficacy of ChatGPT in Therapy Enhancement} &
MDPI(Healthcare) \\
\midrule
46 &
\citealp{Na2024CBT-LLM}. \textit{CBT-LLM: A Chinese Large Language Model for Cognitive Behavioral Therapy-based Mental Health Question Answering} &
ACL \\
\midrule
47 &
\citealp{Sharma2023Cognitive}. \textit{Cognitive Reframing of Negative Thoughts through Human-Language Model Interaction} &
ACL \\
\midrule
48 &
\citealp{Kang2024Can}. \textit{Can Large Language Models be Good Emotional Supporter? Mitigating Preference Bias on Emotional Support Conversation} &
ACL \\
\midrule
49 &
\citealp{Chen2023LLM-empowered}. \textit{LLM-empowered Chatbots for Psychiatrist and Patient Simulation: Application and Evaluation} &
arXiv \\
\midrule
50 &
\citealp{Schmidgall2025AgentClinic}. \textit{AgentClinic: a multimodal agent benchmark to evaluate AI in simulated clinical environments} &
arXiv \\
\midrule
51 &
\citealp{Wang2025Psychological}. \textit{Psychological Counseling CanNot Be Achieved Overnight: Automated Psychological Counseling Through Multi-Session Conversations} &
arXiv \\
\midrule
52 &
\citealp{Tanana2019Development}. \textit{Development and Evaluation of ClientBot: Patient-Like Conversational Agent to Train Basic Counseling Skills} &
JMIR \\
\bottomrule
\end{tabular}
\caption{Papers included in the 52-paper review corpus (Part 3 of 3).}
\label{tab:reviewed-corpus-3}
\end{table*}


\section{Coding Scheme}
\label{app:codebook}

The codebook contains 40 coding fields organized into seven groups:
system and simulation design, evaluation methods, evaluation dimensions,
interaction and behavioral evaluation, theory, transparency and reliability,
additional evaluation characteristics, and claim--evidence alignment. 
Bibliographic metadata and traceability information, such as claim/evidence
quotations and page locations, were recorded separately and are not counted
among the 40 coding fields. 
Tables~\ref{tab:codebook-system}--\ref{tab:codebook-claims} provide the
field definitions, allowable values, and coding rules.

\subsection{System and Simulation Design}

These fields characterize the primary role and design of the simulated system, including what is simulated, how personas or behavioral states are represented, and whether those states change over interaction. 
Table~\ref{tab:codebook-system} summarizes the corresponding fields and coding rules.

\begin{table*}[p]
\centering
\begin{tabular}{p{4.4cm}p{2.2cm}p{8.1cm}}
\toprule
\textbf{Field} & \textbf{Values / Type} & \textbf{Definition and coding rule} \\
\midrule

\texttt{Focus\_Type}
&
Open classification
&
Primary research focus of the paper. Examples in the codebook include Simulation Framework, Evaluation Framework, Dataset, Training System, and Tool.
\\
\midrule

\texttt{Simulation\_Target}
&
Client Agent / Therapist Agent / Dual-Agent / Human Trainee
&
Object or participant role simulated by the system.
\\
\midrule

\texttt{Persona\_Model\_Depth}
&
Surface Persona / Behavioral State Model / Cognitive Model / Dynamic Traits / Expert Principles
&
Depth of the simulated user or client model. Surface Persona uses profile-level descriptions; Behavioral State Model represents behavioral states; Cognitive Model includes beliefs or motivations; Dynamic Traits change over interaction; Expert Principles use expert-defined behavioral principles.
\\
\midrule

\texttt{Uses\_Dynamic\_State}
&
Yes / No
&
Whether the system explicitly tracks and updates a psychological or user state, such as state transitions, session memory, dynamic openness, or related mechanisms.
\\
\midrule

\texttt{Temporal\_Modeling\_Details}
&
Free text
&
Description of how dynamic or temporal state is modeled, such as state transitions, session memory, dynamic openness, or updates to conversational principles.
\\

\bottomrule
\end{tabular}
\caption{Coding fields for system and simulation design.}
\label{tab:codebook-system}
\end{table*}


\subsection{Evaluation Methods}

These fields record who or what performs the evaluation, distinguishing domain experts, lay users, interactive user studies, LLM judges, and automatic metrics. 
Table~\ref{tab:codebook-methods} provides the definitions and coding rules.

\begin{table*}[p]
\centering
\begin{tabular}{p{4.2cm}p{2.2cm}p{8.3cm}}
\toprule
\textbf{Field} & \textbf{Values / Type} & \textbf{Definition and coding rule} \\
\midrule

\texttt{Eval\_Human\_Experts}
&
Yes / No
&
Evaluation by domain experts such as therapists, clinicians, domain specialists, or trained annotators. Excludes crowd workers, participants without domain expertise, and LLM evaluators.
\\
\midrule

\texttt{Eval\_Lay\_Users}
&
Yes / No
&
Evaluation by ordinary users or participants without relevant domain expertise, including crowd workers or participants whose professional expertise is not specified.
\\
\midrule

\texttt{Eval\_User\_Study}
&
Yes / No
&
Participants actively interact with the system under a defined task, protocol, or experimental design. Static output-rating tasks without interaction are excluded.
\\
\midrule

\texttt{Eval\_LLM\_Judge}
&
Yes / No
&
An LLM such as GPT-4 or Claude is prompted to evaluate, score, compare, or rank system outputs. Embedding similarity, feature extraction, or algorithmic metrics alone do not count.
\\
\midrule

\texttt{Eval\_Automatic}
&
Yes / No
&
Evaluation calculated automatically through an algorithm or formula, such as BLEU, ROUGE, Accuracy, F1, cosine similarity, or task success rate. Human and LLM judgments are excluded.
\\

\bottomrule
\end{tabular}
\caption{Coding fields for evaluation methods.}
\label{tab:codebook-methods}
\end{table*}


\subsection{Evaluation Dimensions}

These fields capture what aspects of system performance are explicitly evaluated, including safety, realism, consistency, fidelity, human outcomes, utility, and emotional plausibility. 
Table~\ref{tab:codebook-dimensions} also includes the free-text field used to record reported evaluation metrics.

\begin{table*}[p]
\centering
\begin{tabular}{p{4.3cm}p{1.9cm}p{8.5cm}}
\toprule
\textbf{Field} & \textbf{Values/Type} & \textbf{Definition and coding rule} \\
\midrule

\texttt{Safety}
&
Yes / No
&
Whether the evaluation assesses harmful, biased, unsafe, or ethically problematic outputs. General usefulness alone does not count.
\\
\midrule

\texttt{Realism}
&
Yes / No
&
Whether outputs or behaviors are evaluated as realistic, human-like, natural, or believable. Correctness or task success alone does not count.
\\
\midrule

\texttt{Consistency}
&
Yes / No
&
Whether behavior is evaluated for stability across turns, such as persona consistency, behavioral drift, or trait alignment. Single-turn quality alone does not count.
\\
\midrule

\texttt{Fidelity}
&
Yes / No
&
Whether outputs conform to a predefined identity, profile, role, or target model. General realism alone does not count.
\\
\midrule

\texttt{Human Learning / Outcomes}
&
Yes / No
&
Whether evaluation measures substantive change in human knowledge, skills, or behavior. Perception-only measures are excluded.
\\
\midrule

\texttt{Utility}
&
Yes / No
&
Whether the system is evaluated for usefulness, effectiveness, helpfulness, or task performance. Realism alone does not count.
\\
\midrule

\texttt{Emotional Plausibility}
&
Yes / No
&
Whether emotional behavior or empathic responses are evaluated as realistic or appropriate. Factual correctness alone does not count.
\\
\midrule

\texttt{Raw Eval Metrics}
&
Free text
&
Evaluation metrics reported by the paper, recorded verbatim or in normalized form for descriptive analysis.
\\

\bottomrule
\end{tabular}
\caption{Coding fields for evaluation dimensions and reported metrics.}
\label{tab:codebook-dimensions}
\end{table*}

\subsection{Interaction and Behavioral Evaluation}

These fields characterize how much interaction is represented in the evaluation, how deeply behavior is analyzed across interaction, and whether behavior is tested under different input or intervention conditions. 
Table~\ref{tab:codebook-behavior} defines these fields and their coding categories.

\begin{table*}[p]
\centering
\begin{tabular}{p{4.3cm}p{1.9cm}p{8.5cm}}
\toprule
\textbf{Field} & \textbf{Values/Type} & \textbf{Definition and coding rule} \\
\midrule

\texttt{Interaction\_Level}
&
Single-turn / Short Multi-turn / Extended Dialogue / Longitudinal
&
Amount of interaction represented in the evaluation. Single-turn evaluates one response without dialogue history; Short Multi-turn contains 2--5 turns; Extended Dialogue uses longer conversations, typically 6 or more turns; Longitudinal evaluation spans multiple sessions or a longer period. When multiple interaction levels were present, the highest level demonstrated in the evaluation was coded. 
The six-turn threshold is a pragmatic coding convention used to distinguish short from extended within-session interactions, rather than a theoretical boundary for behavioral change.
\\
\midrule

\texttt{Behavior\_Eval\_Depth}
&
Dynamic / Pattern-level / Static / None
&
Depth of behavioral evaluation.
Dynamic examines how behavior develops across ordered observations within an interaction, such as turn-by-turn trends, transitions between dialogue phases, or systematic qualitative analysis of behavioral progression.
Pattern-level summarizes behavioral frequencies or patterns without tracking change; Static scores behavior independently without examining cross-turn relations; None evaluates output quality without behavioral analysis.
Independent response ratings, aggregate behavioral frequencies, and pre--post endpoint comparisons without analysis of the intervening interaction do not by themselves qualify as Dynamic.
When multiple forms of behavioral evaluation were present, the highest demonstrated evaluation depth was coded.
\\
\midrule

\texttt{Intervention\_Sensitivity}
&
Yes / No
&
Whether the paper explicitly tests differences in system behavior under different inputs or intervention conditions. Papers that do not test or report such behavioral differences are coded No.
\\

\bottomrule
\end{tabular}
\caption{Coding fields for interaction structure and behavioral evaluation.}
\label{tab:codebook-behavior}
\end{table*}


\subsection{Theory, Transparency, and Reliability}

These fields capture theory grounding and operationalization, the types of theory-related evaluation used, prompt disclosure, clinical theory, and reporting of evaluator agreement or reliability. 
Theory grounding captures whether established theory substantively informs the system or study, whereas theory operationalization captures whether that theory is translated into measurable evaluation constructs, criteria, or procedures. 
Table~\ref{tab:codebook-theory} summarizes the corresponding definitions and coding rules.

\begin{table*}[p]
\centering
\begin{tabular}{p{4.3cm}p{2.5cm}p{8.0cm}}
\toprule
\textbf{Field} & \textbf{Values / Type} & \textbf{Definition and coding rule} \\
\midrule

\texttt{Theory\_Operationalized}
&
Strong / Partial / None
&
Degree to which theory is translated into measurable constructs, rubrics, behavioral criteria, or evaluation procedures.
Strong indicates that theory explicitly defines substantial parts of the evaluation; Partial indicates that theory is translated into evaluation criteria only incompletely or for a subset of the theoretically relevant constructs; None indicates that theory may inform the system or study but is not operationalized in the evaluation.
\\
\midrule

\texttt{Prompt\_Disclosure}
&
Full / Partial / No Disclosure
&
Extent to which prompts needed to understand or reproduce the system and evaluation are reported.
Full Disclosure provides the necessary prompts; Partial Disclosure provides incomplete prompts or examples; No Disclosure provides no usable prompt information.
\\
\midrule

\texttt{Theory\_Grounding}
&
Strong / Partial / Mentioned
&
Extent to which established psychological, clinical, counseling, or related theory substantively informs the system or study.
Strong indicates that theory clearly shapes the system design, target constructs, or study rationale; Partial indicates that theory informs the work but only incompletely or indirectly; Mentioned indicates that theory appears mainly as background or motivation without substantively shaping the system or study.
\\
\midrule

\texttt{Clinical\_Theory}
&
Free text
&
Specific psychological, clinical, or counseling theory or framework identified in the paper, such as Cognitive Behavioral Therapy or Motivational Interviewing.
\\
\midrule

\texttt{Reliability\_Reported}
&
Yes / No / N/A
&
Whether evaluator agreement or inter-rater reliability is explicitly reported.
Yes indicates a reported agreement statistic or percentage; No indicates human evaluation without reported agreement; N/A applies when no human evaluation is used.
\\
\midrule

\texttt{Agreement Method}
&
Free text
&
Specific agreement or reliability method reported in the paper, such as Cohen's kappa, Krippendorff's alpha, Fleiss' kappa, or percentage agreement. Left blank when no method is reported.
\\
\midrule

\texttt{Theory\_Eval\_Type}
&
Multi-label categorical
&
Type or types of theory-related evaluation used in the paper.
Categories are not mutually exclusive; the evaluation types relevant to the reported analyses are summarized in Table~\ref{tab:eval-types}.
\\

\bottomrule
\end{tabular}
\caption{Coding fields for theory grounding, theory-related evaluation, transparency, and reliability.}
\label{tab:codebook-theory}
\end{table*}

\subsection{Additional Evaluation Characteristics}

These fields capture additional characteristics of evaluation practice and transparency, including rubric use, LLM-judge validation, comparability, longitudinal and robustness testing, failure analysis, simulation realism, and dataset availability. 
Table~\ref{tab:codebook-quality} provides the corresponding coding rules.

\begin{table*}[p]
\centering
\begin{tabular}{p{4.3cm}p{1.9cm}p{8.5cm}}
\toprule
\textbf{Field} & \textbf{Values/Type} & \textbf{Definition and coding rule} \\
\midrule

\texttt{Has\_Rubric}
&
Yes / No
&
Whether the evaluation provides explicit scoring criteria, a rating scale, rubric, or scoring guidance. Merely stating that annotators provided ratings without criteria is coded No.
\\
\midrule

\texttt{LLM\_Judge\_Validated}
&
Yes / No
&
Whether LLM-judge scores are compared with human judgments in any form, including correlation, agreement tests, mean comparisons, inter-rater tables, or qualitative statements of agreement. Human--human reliability alone does not count. This field records whether any comparison with human judgments was reported and does not imply formal psychometric validation.
\\
\midrule

\texttt{Uses\_Standard\_Metrics}
&
Yes / No
&
Whether established evaluation metrics such as BLEU, ROUGE, or BERTScore are used. Papers using only custom metrics without comparison to standard metrics are coded No.
\\
\midrule

\texttt{Metric\_Interpretable}
&
Yes / No
&
Whether the meaning of the evaluation metric and the interpretation of its scores are clear to the reader.
\\
\midrule

\texttt{Comparable\_To\_Prior\_Work}
&
Yes / No
&
Whether the evaluation uses metrics or datasets that permit comparison with prior work. A custom dataset combined with custom metrics and no baseline is coded No.
\\
\midrule

\texttt{Has\_Longitudinal\_Eval}
&
Yes / No
&
Whether evaluation spans multiple sessions or a longer period. A single-session evaluation is coded No.
\\
\midrule

\texttt{Has\_Robustness\_Testing}
&
Yes / No
&
Whether the system is tested under multiple conditions or datasets rather than only under ideal conditions.
\\
\midrule

\texttt{Has\_Failure\_Analysis}
&
Yes / No
&
Whether the paper explicitly analyzes failure cases, for example in a dedicated failure-analysis section or table.
\\
\midrule

\texttt{Sim\_Behavior\_Realistic}
&
Yes / No
&
Whether simulated behavior includes realistic characteristics such as resistance or emotional fluctuation rather than being overly compliant or mechanical.
\\
\midrule

\texttt{Dataset\_Available}
&
Yes / No
&
Whether the dataset is publicly available through a download link or repository.
\\

\bottomrule
\end{tabular}
\caption{Coding fields for additional evaluation characteristics.}
\label{tab:codebook-quality}
\end{table*}

\subsection{Claim--Evidence Alignment}
\label{app:claim-coding}

The final two fields capture the relationship between the conclusions drawn
in a paper and the evidence reported to support them. Unlike the other
paper-level fields, these fields were applied to predefined claim roles:
the primary empirical or system-level claim (C1) and, where present,
an explicit higher-scope inferential extension beyond C1 (C2).

\begin{table*}[p]
\centering
\small
\begin{tabular}{p{4.2cm}p{3.0cm}p{7.5cm}}
\toprule
\textbf{Field} & \textbf{Values / Type} &
\textbf{Definition and coding rule} \\
\midrule

\texttt{Claim}
&
C1 / C2; free text
&
Selected empirical claim and its inferential role.
C1 captures the paper's primary empirical or system-level conclusion.
C2 is optional and captures the strongest explicit, non-speculative
inferential extension beyond C1.
Motivational statements, design goals, hypotheses, recommendations,
future work, and purely speculative statements are excluded.
\\
\midrule

\texttt{Claim\_Evidence\_Alignment}
&
Aligned / Partial / Exceeds / Unclear
&
Degree to which the strongest reported evidence supports the selected claim.
\textsc{Aligned} indicates that the evidence directly evaluates the claimed
target at the claimed scope and evidential level.
\textsc{Partial} indicates directly relevant evidence for the same fundamental
target, but with narrower scope, population, duration, or outcome.
\textsc{Exceeds} indicates that the claim extends to a materially different
construct, population, outcome, or evidential level.
\textsc{Unclear} indicates insufficient or contradictory reporting for a
reliable judgment.
\\

\bottomrule
\end{tabular}
\caption{Coding fields for claim--evidence alignment.}
\label{tab:codebook-claims}
\end{table*}

Claims were selected primarily from the abstract, explicit contribution
statements, and conclusion; discussion sections were used only when they
contained an explicit, nonredundant empirical conclusion. 
C2 was identified when a paper drew an explicit conclusion extending
beyond its primary empirical conclusion (C1), regardless of whether
the available evidence supported that extension. 
We distinguished such conclusions from speculative future applications:
statements describing potential uses, research objectives, or outcomes
requiring future investigation were excluded. 
Hedged language alone did not determine eligibility; the key criterion
was whether the statement was presented as an inference from the
reported findings rather than a prospective possibility. 
Importantly, a C2 claim could be classified as \textsc{Aligned} if
the reported evidence directly supported its broader conclusion.

For each selected claim, the full paper was examined to identify
the strongest directly relevant supporting evidence. 
Alignment judgments followed an evidence-level preservation principle:
evidence at the response, simulation, perception, behavioral, or
human-outcome level was treated as support for claims at that evidential
level, rather than automatically as support for stronger downstream
conclusions. 
Claim and evidence quotations and their page or section locations
were retained for traceability.

Operationally, we distinguished \textsc{Partial} from \textsc{Exceeds}
by asking whether the claim could be narrowed to match the reported
evidence without changing its fundamental construct, outcome, population,
or evidential level. If so, it was coded \textsc{Partial}; if narrowing
would change the fundamental target of the claim, it was coded
\textsc{Exceeds}.

Table~\ref{tab:claim-evidence-examples} illustrates how these rules were
applied to claims in the reviewed corpus.

\begin{table*}[p]
\centering
\small
\begin{tabular}{p{2.4cm}p{4.2cm}p{5.8cm}p{1.9cm}}
\toprule
\textbf{Paper / role} &
\textbf{Selected claim} &
\textbf{Strongest supporting evidence} &
\textbf{Alignment} \\
\midrule

\citealp{Liu2025Eeyore}, C1
&
Eeyore outperforms GPT-4o-based simulation baselines in linguistic
authenticity and profile adherence.
&
Fifteen counselors or advanced psychology students evaluated live
interactions on unseen profiles for authenticity and profile adherence
(Sec.~4.2; Fig.~4). The evaluated outcomes therefore directly correspond
to the comparative claim.
&
\textsc{Aligned}
\\
\midrule

\citealp{Sungarda2025GenAI}, C2
&
SWITCH provides a scalable, low-cost, and consistent training workflow
that can complement field education.
&
The system was introduced in a compulsory social-work practice course
and teacher workshops, supporting its use as a complementary training
workflow; however, cost, scalability, and consistency were not directly
measured (Sec.~VI; Tables~VI--VII).
&
\textsc{Partial}
\\
\midrule

\citealp{Chen2025MIND}, C2
&
MIND's inner-dialogue interactions enhance empathy and emotional resonance,
reduce self-criticism, and strengthen self-identity.
&
A short-term study with eight participants measured affect using PANAS
and subjective user-experience ratings, together with expert ratings of
generated content; self-criticism and self-identity were not directly
measured (Sec.~3.4; Tables~3--4).
&
\textsc{Exceeds}
\\

\bottomrule
\end{tabular}
\caption{Illustrative applications of the claim--evidence alignment coding.
The examples show how alignment depends on whether the reported evidence
preserves the target and evidential scope of the selected claim. They are
illustrative rather than exhaustive; complete claim--evidence records are
included in the released coding materials.}
\label{tab:claim-evidence-examples}
\end{table*}

\section{Coding Reliability}
\label{app:reliability}

For the 40 fields, two researchers independently coded
the full corpus before adjudication. Mean raw agreement across fields was
82.5\%; among fields for which Cohen's $\kappa$ was computable, mean
$\kappa$ was 0.61. Disagreements and ambiguous cases were resolved through
discussion with the faculty lead.

\section{Supplementary Results}
\label{app:supp-results}

% YOUR EXISTING appendices.tex begins here
\begin{table*}[tbp]
\centering
\small
\begin{tabular}{lc}
\toprule
\textbf{Evaluator Type} & \textbf{Count (\%)} \\
\midrule
Human Experts & 34 (65.4\%) \\
LLM Judges & 30 (57.7\%) \\
Automatic Metrics & 35 (67.3\%) \\
All Three (Expert + LLM + Auto) & 15 (28.8\%) \\
\bottomrule
\end{tabular}
\caption{Distribution of evaluator types across reviewed studies ($N=52$). Counts indicate the number of papers using each evaluator type; categories are not mutually exclusive, except for the final row, which reports papers combining human experts, LLM judges, and automatic metrics.}\label{tab:evaluator-types}
\end{table*}

\begin{table*}[tbp]
\centering
\small
\begin{tabular}{p{5cm}p{10cm}}
\toprule
\textbf{Evaluator Type} & \textbf{Papers} \\
\midrule
Human Experts & \citealp{Aghakhani2025From}, \citealp{Berrezueta-Guzman2024Future}, \citealp{Cai2026Exploring}, \citealp{Chen2023LLM-empowered}, \citealp{Chen2025MIND}, \citealp{Deshpande2025Evaluating}, \citealp{Feng2025Reframe}, \citealp{Gabriel2024Can}, \citealp{Haydarov2025Towards}, \citealp{Kim2025Can}, \citealp{Kim2025KMI}, \citealp{Kim2025MIRROR}, \citealp{Kuhail2025Human-Human}, \citealp{Lee2024CACTUS}, \citealp{Li2024Understanding}, \citealp{Liu2025Eeyore}, \citealp{Louie2024Roleplay-doh}, \citealp{Louie2026Can}, \citealp{Na2024CBT-LLM}, \citealp{Nguyen2025Do}, \citealp{Qiu2026PsyCLIENT}, \citealp{Samuel2025PersonaGym}, \citealp{Schmidgall2025AgentClinic}, \citealp{Sharma2023Cognitive}, \citealp{Wang2024PATIENT}, \citealp{Wang2025CARE-Bench}, \citealp{Wang2025Psychological}, \citealp{Xiao2024HealMe}, \citealp{Xie2025PsyDT}, \citealp{Xu2025MentalChat16K}, \citealp{Xu2025MultiAgentESC}, \citealp{Yang2025CAMI}, \citealp{Yang2025Consistent}, \citealp{Zhang2025CBT-BENCH} \\
\midrule
LLM Judges & \citealp{Aghakhani2025From}, \citealp{Cai2026Exploring}, \citealp{Chiu2024A}, \citealp{Feng2025Reframe}, \citealp{Fouda2026PsychiatryBench}, \citealp{Haydarov2025Towards}, \citealp{Kim2025Can}, \citealp{Kim2025MIRROR}, \citealp{Lee2024CACTUS}, \citealp{Li2024Understanding}, \citealp{Liu2023ChatCounselor}, \citealp{Liu2025Eeyore}, \citealp{Louie2026Can}, \citealp{Nguyen2025Do}, \citealp{Ozgun2025Trustworthy}, \citealp{Qiu2025EmoAgent}, \citealp{Qiu2026PsyCLIENT}, \citealp{Samuel2025PersonaGym}, \citealp{Schmidgall2025AgentClinic}, \citealp{Wang2024PATIENT}, \citealp{Wang2024Towards}, \citealp{Wang2025AnnaAgent}, \citealp{Wang2025CARE-Bench}, \citealp{Wang2025Psychological}, \citealp{Xiao2024HealMe}, \citealp{Xie2025PsyDT}, \citealp{Xu2025MentalChat16K}, \citealp{Xu2025MultiAgentESC}, \citealp{Yang2025CAMI}, \citealp{Yang2025Consistent} \\
\midrule
Automatic Metrics & \citealp{Aghakhani2025From}, \citealp{Basar2025How}, \citealp{Chen2023LLM-empowered}, \citealp{Chiu2024A}, \citealp{Deshpande2025Evaluating}, \citealp{Feng2025Reframe}, \citealp{Fouda2026PsychiatryBench}, \citealp{Gabriel2024Can}, \citealp{Hu2024PsycoLLM}, \citealp{Kang2024Can}, \citealp{Kermani2025Systematic}, \citealp{Kim2025Can}, \citealp{Li2024Understanding}, \citealp{Louie2026Can}, \citealp{Na2024CBT-LLM}, \citealp{Nguyen2025Do}, \citealp{Ozgun2025Trustworthy}, \citealp{Qiu2025EmoAgent}, \citealp{Samuel2025PersonaGym}, \citealp{Schmidgall2025AgentClinic}, \citealp{Schmidmaier2025UsingNonverbal}, \citealp{Sharma2023Cognitive}, \citealp{Sungarda2025GenAI}, \citealp{Tanana2019Development}, \citealp{Wang2024PATIENT}, \citealp{Wang2024Towards}, \citealp{Wang2025AnnaAgent}, \citealp{Wang2025CARE-Bench}, \citealp{Wang2025Feel}, \citealp{Wang2025Psychological}, \citealp{Xie2025PsyDT}, \citealp{Xu2025MultiAgentESC}, \citealp{Yang2025CAMI}, \citealp{Yang2025Consistent}, \citealp{Zhang2025CBT-BENCH} \\
\midrule
All Three (Expert + LLM + Auto) & \citealp{Aghakhani2025From}, \citealp{Feng2025Reframe}, \citealp{Kim2025Can}, \citealp{Li2024Understanding}, \citealp{Louie2026Can}, \citealp{Nguyen2025Do}, \citealp{Samuel2025PersonaGym}, \citealp{Schmidgall2025AgentClinic}, \citealp{Wang2024PATIENT}, \citealp{Wang2025CARE-Bench}, \citealp{Wang2025Psychological}, \citealp{Xie2025PsyDT}, \citealp{Xu2025MultiAgentESC}, \citealp{Yang2025CAMI}, \citealp{Yang2025Consistent} \\
\bottomrule
\end{tabular}
\caption{Evaluator types used in reviewed LLM-based counseling and role-play simulation papers. Papers are grouped by whether evaluation involved human experts, LLM judges, automatic metrics, or all three evaluator types.}
\label{tab:evaluator-types-cite}
\end{table*}
%%%%%%%%%%%%%%authority


\begin{table*}[tbp]
\centering
\small
\begin{tabular}{lc}
\toprule
\textbf{Evaluation Dimension} & \textbf{Count (\%)} \\
\midrule
Utility & 44 (84.6\%) \\
Fidelity & 33 (63.5\%) \\
Realism & 29 (55.8\%) \\
Emotional Plausibility & 28 (53.8\%) \\
Consistency & 20 (38.5\%) \\
Safety & 14 (26.9\%) \\
Human Learning / Outcomes & 10 (19.2\%) \\
\bottomrule
\end{tabular}
\caption{Evaluation dimensions reported across reviewed studies ($N=52$). Utility, fidelity, realism, and emotional plausibility were more commonly evaluated than consistency, safety, and human learning or outcome measures; categories are not mutually exclusive.}
\label{tab:eval-dimensions}
\end{table*}

\begin{table*}[tbp]
\centering
\small
\begin{tabular}{p{4cm}p{10.5cm}}
\toprule
\textbf{Category} & \textbf{Papers} \\
\midrule
Utility & \citealp{Aghakhani2025From}, \citealp{Basar2024To}, \citealp{Basar2025How}, \citealp{Berrezueta-Guzman2024Future}, \citealp{Cai2026Exploring}, \citealp{Chen2023LLM-empowered}, \citealp{Chen2025MIND}, \citealp{Deshpande2025Evaluating}, \citealp{Feng2025Reframe}, \citealp{Fouda2026PsychiatryBench}, \citealp{Gabriel2024Can}, \citealp{Haydarov2025Towards}, \citealp{Hu2024PsycoLLM}, \citealp{Jaiswal2024Building}, \citealp{Kang2024Can}, \citealp{Kermani2025Systematic}, \citealp{Kim2025Can}, \citealp{Kim2025KMI}, \citealp{Kim2025MIRROR}, \citealp{Kuhail2025Human-Human}, \citealp{Lee2024CACTUS}, \citealp{Li2024Understanding}, \citealp{Liu2023ChatCounselor}, \citealp{Louie2024Roleplay-doh}, \citealp{Louie2026Can}, \citealp{Na2024CBT-LLM}, \citealp{Nguyen2025Do}, \citealp{Ozgun2025Trustworthy}, \citealp{Qiu2025EmoAgent}, \citealp{Qiu2026PsyCLIENT}, \citealp{Schmidgall2025AgentClinic}, \citealp{Schmidmaier2025Using}, \citealp{Sharma2023Cognitive}, \citealp{Tanana2019Development}, \citealp{Wang2024PATIENT}, \citealp{Wang2024Towards}, \citealp{Wang2025CARE-Bench}, \citealp{Wang2025Psychological}, \citealp{Xiao2024HealMe}, \citealp{Xie2025PsyDT}, \citealp{Xu2025MentalChat16K}, \citealp{Xu2025MultiAgentESC}, \citealp{Yang2025CAMI}, \citealp{Zhang2025CBT-BENCH} \\
\midrule
Fidelity & \citealp{Chen2023LLM-empowered}, \citealp{Chen2025MIND}, \citealp{Chiu2024A}, \citealp{Deshpande2025Evaluating}, \citealp{Feng2025Reframe}, \citealp{Fouda2026PsychiatryBench}, \citealp{Gabriel2024Can}, \citealp{Kang2024Can}, \citealp{Kim2025Can}, \citealp{Kim2025KMI}, \citealp{Kim2025MIRROR}, \citealp{Li2024Understanding}, \citealp{Liu2023ChatCounselor}, \citealp{Liu2025Eeyore}, \citealp{Louie2024Roleplay-doh}, \citealp{Louie2026Can}, \citealp{Na2024CBT-LLM}, \citealp{Ozgun2025Trustworthy}, \citealp{Qiu2026PsyCLIENT}, \citealp{Samuel2025PersonaGym}, \citealp{Schmidgall2025AgentClinic}, \citealp{Tanana2019Development}, \citealp{Wang2024PATIENT}, \citealp{Wang2024Towards}, \citealp{Wang2025AnnaAgent}, \citealp{Wang2025CARE-Bench}, \citealp{Wang2025Feel}, \citealp{Wang2025Psychological}, \citealp{Xiao2024HealMe}, \citealp{Xie2025PsyDT}, \citealp{Yang2025CAMI}, \citealp{Yang2025Consistent}, \citealp{Zhang2025CBT-BENCH} \\
\midrule
Realism & \citealp{Basar2024To}, \citealp{Basar2025How}, \citealp{Berrezueta-Guzman2024Future}, \citealp{Chen2023LLM-empowered}, \citealp{Chen2025MIND}, \citealp{Chiu2024A}, \citealp{Deshpande2025Evaluating}, \citealp{Kim2025Can}, \citealp{Kim2025KMI}, \citealp{Kuhail2025Human-Human}, \citealp{Lee2024CACTUS}, \citealp{Liu2023ChatCounselor}, \citealp{Liu2025Eeyore}, \citealp{Louie2024Roleplay-doh}, \citealp{Louie2026Can}, \citealp{Ozgun2025Trustworthy}, \citealp{Qiu2026PsyCLIENT}, \citealp{Samuel2025PersonaGym}, \citealp{Schmidgall2025AgentClinic}, \citealp{Tanana2019Development}, \citealp{Wang2024PATIENT}, \citealp{Wang2024Towards}, \citealp{Wang2025AnnaAgent}, \citealp{Wang2025CARE-Bench}, \citealp{Wang2025Feel}, \citealp{Xie2025PsyDT}, \citealp{Xu2025MultiAgentESC}, \citealp{Yang2025CAMI}, \citealp{Yang2025Consistent} \\
\midrule
Emotional Plausibility & \citealp{Berrezueta-Guzman2024Future}, \citealp{Cai2026Exploring}, \citealp{Chen2023LLM-empowered}, \citealp{Chen2025MIND}, \citealp{Gabriel2024Can}, \citealp{Haydarov2025Towards}, \citealp{Kang2024Can}, \citealp{Kim2025Can}, \citealp{Kim2025MIRROR}, \citealp{Kuhail2025Human-Human}, \citealp{Liu2023ChatCounselor}, \citealp{Liu2025Eeyore}, \citealp{Ozgun2025Trustworthy}, \citealp{Qiu2026PsyCLIENT}, \citealp{Schmidgall2025AgentClinic}, \citealp{Schmidmaier2025Using}, \citealp{Sharma2023Cognitive}, \citealp{Wang2024PATIENT}, \citealp{Wang2024Towards}, \citealp{Wang2025CARE-Bench}, \citealp{Wang2025Feel}, \citealp{Wang2025Psychological}, \citealp{Xiao2024HealMe}, \citealp{Xie2025PsyDT}, \citealp{Xu2025MentalChat16K}, \citealp{Xu2025MultiAgentESC}, \citealp{Yang2025CAMI}, \citealp{Zhang2025CBT-BENCH} \\
\midrule
Consistency & \citealp{Berrezueta-Guzman2024Future}, \citealp{Chen2023LLM-empowered}, \citealp{Chen2025MIND}, \citealp{Chiu2024A}, \citealp{Fouda2026PsychiatryBench}, \citealp{Gabriel2024Can}, \citealp{Haydarov2025Towards}, \citealp{Kim2025KMI}, \citealp{Li2024Understanding}, \citealp{Liu2025Eeyore}, \citealp{Louie2024Roleplay-doh}, \citealp{Ozgun2025Trustworthy}, \citealp{Qiu2026PsyCLIENT}, \citealp{Wang2024Towards}, \citealp{Wang2025AnnaAgent}, \citealp{Wang2025CARE-Bench}, \citealp{Wang2025Psychological}, \citealp{Xiao2024HealMe}, \citealp{Yang2025CAMI}, \citealp{Yang2025Consistent} \\
\midrule
Safety & \citealp{Basar2024To}, \citealp{Basar2025How}, \citealp{Berrezueta-Guzman2024Future}, \citealp{Cai2026Exploring}, \citealp{Fouda2026PsychiatryBench}, \citealp{Gabriel2024Can}, \citealp{Li2024Understanding}, \citealp{Qiu2025EmoAgent}, \citealp{Samuel2025PersonaGym}, \citealp{Sharma2023Cognitive}, \citealp{Wang2024Towards}, \citealp{Xiao2024HealMe}, \citealp{Xie2025PsyDT}, \citealp{Xu2025MentalChat16K} \\
\midrule
Human Learning / Outcomes & \citealp{Berrezueta-Guzman2024Future}, \citealp{Chen2025MIND}, \citealp{Feng2025Reframe}, \citealp{Li2024Understanding}, \citealp{Louie2026Can}, \citealp{Qiu2026PsyCLIENT}, \citealp{Tanana2019Development}, \citealp{Wang2024PATIENT}, \citealp{Wang2025Psychological}, \citealp{Xiao2024HealMe} \\
\bottomrule
\end{tabular}
\caption{Reviewed papers by evaluation dimension. Papers may appear in multiple rows because evaluation dimensions are not mutually exclusive.}
\label{tab:eval-dimensions-cite}
\end{table*}
%%%%%%%%%%local outputs
% === Theory Grounding & Operationalization Tables ===

\begin{table*}[tbp]
\centering
\begin{tabular}{lccccc}
\toprule
\textbf{Category} &
\textbf{Strong} &
\textbf{Partial} &
\textbf{Mentioned} &
\textbf{None} &
\textbf{Total} \\
\midrule

Theory Grounding &
40 (76.9\%) &
10 (19.2\%) &
2 (3.8\%) &
--- &
52 \\

\midrule

Theory Operationalization &
30 (57.7\%) &
16 (30.8\%) &
--- &
6 (11.5\%) &
52 \\

\bottomrule
\end{tabular}
\caption{Distribution of theory grounding and theory operationalization across reviewed studies. Theory grounding captures whether established theory substantively informs the system or study, whereas theory operationalization captures whether theory is translated into measurable constructs, rubrics, behavioral criteria, or evaluation procedures.}
\label{tab:theory-stats}
\end{table*}

\begin{table*}[tbp]
\centering
\small
\begin{tabular}{p{2.5cm}p{12.4cm}}
\toprule
\textbf{Category} & \textbf{Papers} \\
\midrule
Strong & \citealp{Xiao2024HealMe}, \citealp{Gabriel2024Can}, \citealp{Haydarov2025Towards}, \citealp{Qiu2025EmoAgent}, \citealp{Aghakhani2025From}, \citealp{Wang2025CARE-Bench}, \citealp{Sungarda2025GenAI}, \citealp{Chiu2024A}, \citealp{Li2024Understanding}, \citealp{Wang2024Towards}, \citealp{Fouda2026PsychiatryBench}, \citealp{Cai2026Exploring}, \citealp{Nguyen2025Do}, \citealp{Zhang2025CBT-BENCH}, \citealp{Kim2025KMI}, \citealp{Jaiswal2024Building}, \citealp{Xu2025MentalChat16K}, \citealp{Schmidmaier2025Using}, \citealp{Schmidmaier2025UsingNonverbal}, \citealp{Ozgun2025Trustworthy}, \citealp{Louie2026Can}, \citealp{Wang2024PATIENT}, \citealp{Lee2024CACTUS}, \citealp{Basar2024To}, \citealp{Yang2025Consistent}, \citealp{Yang2025CAMI}, \citealp{Xie2025PsyDT}, \citealp{Basar2025How}, \citealp{Feng2025Reframe}, \citealp{Xu2025MultiAgentESC}, \citealp{Kim2025MIRROR}, \citealp{Kim2025Can}, \citealp{Liu2025Eeyore}, \citealp{Wang2025Feel}, \citealp{Samuel2025PersonaGym}, \citealp{Chen2025MIND}, \citealp{Na2024CBT-LLM}, \citealp{Sharma2023Cognitive}, \citealp{Wang2025Psychological}, \citealp{Tanana2019Development} \\
\midrule
Partial & \citealp{Wang2025AnnaAgent}, \citealp{Liu2023ChatCounselor}, \citealp{Hu2024PsycoLLM}, \citealp{Qiu2026PsyCLIENT}, \citealp{Deshpande2025Evaluating}, \citealp{Louie2024Roleplay-doh}, \citealp{Kermani2025Systematic}, \citealp{Kang2024Can}, \citealp{Chen2023LLM-empowered}, \citealp{Schmidgall2025AgentClinic} \\
\midrule
Mentioned & \citealp{Kuhail2025Human-Human}, \citealp{Berrezueta-Guzman2024Future} \\
\bottomrule
\end{tabular}
\caption{Reviewed papers classified by theory grounding level. Papers are grouped according to whether theory grounding was coded as strong, partial, or mentioned only.}
\label{tab:theory-grounding}
\end{table*}

\begin{table*}[tbp]
\centering
\small
\begin{tabular}{p{2.5cm}p{12.4cm}}
\toprule
\textbf{Category} & \textbf{Papers} \\
\midrule
Strong & \citealp{Xiao2024HealMe}, \citealp{Gabriel2024Can}, \citealp{Qiu2025EmoAgent}, \citealp{Aghakhani2025From}, \citealp{Wang2025CARE-Bench}, \citealp{Sungarda2025GenAI}, \citealp{Chiu2024A}, \citealp{Li2024Understanding}, \citealp{Wang2024Towards}, \citealp{Cai2026Exploring}, \citealp{Nguyen2025Do}, \citealp{Zhang2025CBT-BENCH}, \citealp{Kim2025KMI}, \citealp{Xu2025MentalChat16K}, \citealp{Schmidmaier2025Using}, \citealp{Schmidmaier2025UsingNonverbal}, \citealp{Louie2026Can}, \citealp{Wang2024PATIENT}, \citealp{Lee2024CACTUS}, \citealp{Yang2025CAMI}, \citealp{Feng2025Reframe}, \citealp{Kim2025MIRROR}, \citealp{Kim2025Can}, \citealp{Wang2025Feel}, \citealp{Samuel2025PersonaGym}, \citealp{Na2024CBT-LLM}, \citealp{Sharma2023Cognitive}, \citealp{Kang2024Can}, \citealp{Wang2025Psychological}, \citealp{Tanana2019Development} \\
\midrule
Partial & \citealp{Haydarov2025Towards}, \citealp{Liu2023ChatCounselor}, \citealp{Hu2024PsycoLLM}, \citealp{Fouda2026PsychiatryBench}, \citealp{Qiu2026PsyCLIENT}, \citealp{Jaiswal2024Building}, \citealp{Ozgun2025Trustworthy}, \citealp{Louie2024Roleplay-doh}, \citealp{Basar2024To}, \citealp{Yang2025Consistent}, \citealp{Kermani2025Systematic}, \citealp{Basar2025How}, \citealp{Liu2025Eeyore}, \citealp{Chen2025MIND}, \citealp{Chen2023LLM-empowered}, \citealp{Schmidgall2025AgentClinic} \\
\midrule
None & \citealp{Wang2025AnnaAgent}, \citealp{Deshpande2025Evaluating}, \citealp{Kuhail2025Human-Human}, \citealp{Xie2025PsyDT}, \citealp{Xu2025MultiAgentESC}, \citealp{Berrezueta-Guzman2024Future} \\
\bottomrule
\end{tabular}
\caption{Reviewed papers grouped by theory operationalization level. Operationalization refers to whether theory is translated into measurable constructs, rubrics, behavioral criteria, or evaluation procedures.}
\label{tab:theory-operationalization}
\end{table*}

\begin{table*}[tbp]
\centering
\small
\begin{tabular}{p{6cm}p{8cm}}
\toprule
\textbf{Evaluation Type} & \textbf{Papers} \\
\midrule
Validated scale\newline\textcolor{gray}{\footnotesize{\textit{Standardized psychological or clinical scales}}} & \citealp{Xiao2024HealMe}, \citealp{Qiu2025EmoAgent}, \citealp{Li2024Understanding}, \citealp{Wang2024Towards}, \citealp{Schmidmaier2025Using}, \citealp{Schmidmaier2025UsingNonverbal}, \citealp{Yang2025CAMI}, \citealp{Chen2025MIND}, \citealp{Wang2025Psychological} \\
\midrule
Established therapy/counseling coding system\newline\textcolor{gray}{\footnotesize{\textit{Therapy process or outcome coding}}} & \citealp{Xiao2024HealMe}, \citealp{Zhang2025CBT-BENCH}, \citealp{Lee2024CACTUS}, \citealp{Feng2025Reframe}, \citealp{Kim2025MIRROR}, \citealp{Kim2025Can}, \citealp{Tanana2019Development} \\
\midrule
Theory-specific rubric\newline\textcolor{gray}{\footnotesize{\textit{Custom rubrics based on theoretical constructs}}} & \citealp{Xiao2024HealMe}, \citealp{Gabriel2024Can}, \citealp{Qiu2025EmoAgent}, \citealp{Aghakhani2025From}, \citealp{Wang2025CARE-Bench}, \citealp{Sungarda2025GenAI}, \citealp{Liu2023ChatCounselor}, \citealp{Chiu2024A}, \citealp{Li2024Understanding}, \citealp{Wang2024Towards}, \citealp{Fouda2026PsychiatryBench}, \citealp{Cai2026Exploring}, \citealp{Qiu2026PsyCLIENT}, \citealp{Nguyen2025Do}, \citealp{Zhang2025CBT-BENCH}, \citealp{Kim2025KMI}, \citealp{Xu2025MentalChat16K}, \citealp{Ozgun2025Trustworthy}, \citealp{Louie2026Can}, \citealp{Louie2024Roleplay-doh}, \citealp{Wang2024PATIENT}, \citealp{Lee2024CACTUS}, \citealp{Basar2024To}, \citealp{Yang2025Consistent}, \citealp{Yang2025CAMI}, \citealp{Basar2025How}, \citealp{Feng2025Reframe}, \citealp{Kim2025MIRROR}, \citealp{Kim2025Can}, \citealp{Samuel2025PersonaGym}, \citealp{Chen2025MIND}, \citealp{Na2024CBT-LLM}, \citealp{Sharma2023Cognitive}, \citealp{Kang2024Can}, \citealp{Wang2025Psychological} \\
\midrule
Custom expert rating\newline\textcolor{gray}{\footnotesize{\textit{Expert judgments on custom dimensions}}} & \citealp{Xiao2024HealMe}, \citealp{Gabriel2024Can}, \citealp{Haydarov2025Towards}, \citealp{Wang2025CARE-Bench}, \citealp{Qiu2026PsyCLIENT}, \citealp{Kim2025KMI}, \citealp{Louie2026Can}, \citealp{Louie2024Roleplay-doh}, \citealp{Wang2024PATIENT}, \citealp{Basar2024To}, \citealp{Yang2025Consistent}, \citealp{Basar2025How}, \citealp{Kim2025MIRROR}, \citealp{Liu2025Eeyore}, \citealp{Na2024CBT-LLM}, \citealp{Sharma2023Cognitive}, \citealp{Kang2024Can}, \citealp{Schmidgall2025AgentClinic}, \citealp{Wang2025Psychological} \\
\midrule
Affective/emotional dynamics metric\newline\textcolor{gray}{\footnotesize{\textit{Emotion recognition or affective computing metrics}}} & \citealp{Haydarov2025Towards}, \citealp{Schmidmaier2025UsingNonverbal}, \citealp{Wang2025Feel} \\
\midrule
Clinical diagnostic benchmark\newline\textcolor{gray}{\footnotesize{\textit{Clinical diagnosis or screening benchmarks}}} & \citealp{Qiu2025EmoAgent}, \citealp{Hu2024PsycoLLM}, \citealp{Fouda2026PsychiatryBench}, \citealp{Cai2026Exploring}, \citealp{Ozgun2025Trustworthy}, \citealp{Liu2025Eeyore}, \citealp{Chen2023LLM-empowered}, \citealp{Schmidgall2025AgentClinic} \\
\midrule
Theory-specific benchmark/task\newline\textcolor{gray}{\footnotesize{\textit{Benchmark or task whose scoring criteria or target behaviors are derived from a psychological, clinical, or counseling construct.}}} & \citealp{Qiu2025EmoAgent}, \citealp{Aghakhani2025From}, \citealp{Wang2025CARE-Bench}, \citealp{Sungarda2025GenAI}, \citealp{Chiu2024A}, \citealp{Hu2024PsycoLLM}, \citealp{Fouda2026PsychiatryBench}, \citealp{Nguyen2025Do}, \citealp{Zhang2025CBT-BENCH}, \citealp{Kim2025KMI}, \citealp{Jaiswal2024Building}, \citealp{Xu2025MentalChat16K}, \citealp{Ozgun2025Trustworthy}, \citealp{Wang2024PATIENT}, \citealp{Lee2024CACTUS}, \citealp{Kermani2025Systematic}, \citealp{Feng2025Reframe}, \citealp{Kim2025Can}, \citealp{Liu2025Eeyore}, \citealp{Wang2025Feel}, \citealp{Samuel2025PersonaGym}, \citealp{Na2024CBT-LLM}, \citealp{Sharma2023Cognitive}, \citealp{Kang2024Can}, \citealp{Chen2023LLM-empowered}, \citealp{Tanana2019Development} \\
\bottomrule
\end{tabular}
\caption{Theory-related evaluation types identified across the reviewed studies. Papers may appear in multiple rows because categories are not mutually exclusive.}
\label{tab:eval-types}
\end{table*}


% === Evaluation Methodology Tables ===








